In a companion manuscript on the regular tetrahedron~\cite{BabanskyyTetraCoxeterPure2026}, the $A_3$
Coxeter dissection is controlled by a small but rigid combinatorial
object: posets on four elements.  In that setting, convex chamber-unions
are poset cones, the chamber count is the number of linear extensions, and
the admissible convex tiling pieces are exactly the series-parallel poset
types.  The purpose of this note is to extract the type-$A$ mechanism from
that rank-three situation and ask what remains true in $S_n$.
We use standard terminology for finite posets and linear extensions as in
Trotter~\cite{Trotter1992} and Stanley~\cite{Stanley2012}.

The general dictionary is simple.  In the braid arrangement of type
$A_{n-1}$, choosing a side of the reflecting hyperplane $x_i=x_j$ is the
same as imposing one precedence constraint, either $i<j$ or $j<i$.
Compatible systems of such choices close transitively to partial orders.
The chambers satisfying those constraints are exactly the linear
extensions of the resulting poset.  Thus the halfspace-convex chamber
unions in type $A$ are naturally the poset cones $C(P)$.

The tiling question is more delicate.  If $P$ is a poset on $[n]$, write
$L(P)\subseteq S_n$ for its linear extensions.  We say that $P$ is an
extension tiler when
\[
        S_n=\bigsqcup_{a\in A} aL(P)
\]
for some set $A\subseteq S_n$.  This is a group-translate condition inside
the Coxeter complex.  It should not be confused with a classification of
all possible Euclidean congruent copies in an ambient geometric problem.
The setting is adjacent to earlier work on quotients and splittings in
Coxeter groups~\cite{BjornerWachs1988,GaetzGao2020}, but differs from it
because the complement set $A$ is allowed to be an arbitrary set of left
translators rather than a canonical weak-order quotient.

The first main point is positive and general: every finite
series-parallel poset is an extension tiler.  This is the mechanism already
used in the $A_3$ Coxeter-pure paper; here we state it as a standalone
type-$A$ theorem.

The second main point begins at $n=5$.  In rank four, divisibility stops
being a sufficient proxy for series-parallelity.  There are $51$
unlabeled poset types on five elements for which $|L(P)|$ divides $120$,
but only $48$ of them are series-parallel.  We prove that the remaining
three types do not tile $S_5$ by left translates using elementary
fiber-counting obstructions.  The former rank and packing certificates
are retained only as an independent computational audit.

The resulting general theorem is only one-sided:
\[
\text{series-parallel} \Longrightarrow \text{extension tiler}
\]
for all finite posets.  The converse is known in the previously analyzed
$A_3/S_4$ Coxeter-pure setting and is proved here for $S_5$:
\[
\text{series-parallel} \Longleftrightarrow \text{extension tiler}
\]
in those finite test cases.  However, the next rank already produces a
counterexample: we give a six-element non-series-parallel poset whose
linear extensions tile $S_6$.

We keep three claim levels separate throughout the paper.  The type-$A$
poset-cone dictionary and the series-parallel tiler theorem are ordinary
theorems.  The $S_5$ classification is also an ordinary finite theorem,
with the computer relegated to audit status.  The $S_6$ example is a
finite certificate that the naive converse is false; the remaining
problem is to classify the larger class of extension tilers.
